%!TEX root = tfg.tex
\chapter{Manual de usuario}

\section{Alcance del manual}
El objetivo de este manual es explicar de forma práctica cómo iniciar una partida, cómo
interactuar con la interfaz y con los sistemas principales del juego y qué funciones quedan
disponibles en la versión actual. No se centra en la instalación del proyecto ni en su
configuración técnica, sino en el uso directo del producto una vez ejecutado.

\section{Inicio de la aplicación}
Al arrancar el juego se muestra el menú principal. Desde este punto pueden utilizarse los
botones disponibles en pantalla para acceder a las funciones iniciales del sistema:

\begin{itemize}
    \item \textbf{Iniciar partida:} carga la escena principal y comienza una nueva \textit{run} en modo normal.
    \item \textbf{Dev Mode:} carga la misma escena jugable, pero activa el modo de depuración para habilitar atajos adicionales de desarrollo.
    \item \textbf{Salir:} cierra la aplicación. En el entorno de Unity detiene el modo de juego del editor.
\end{itemize}

\section{Objetivo general de la partida}
La partida se estructura como una exploración progresiva por círculos del Infierno. El
jugador debe recorrer las habitaciones generadas proceduralmente, superar encuentros con
enemigos, recoger recursos y objetos, localizar la llave del círculo actual y alcanzar la
escalera que permite descender al siguiente nivel temático. Este flujo se repite hasta
completar el último círculo, momento en el que se activa la condición de victoria.

Durante la exploración aparecen además salas especiales, como tienda, tesoro o sala secreta,
que introducen recompensas, compras y desvíos dentro del recorrido principal.

\section{Controles principales}
\subsection{Controles con teclado}
La versión jugable se apoya principalmente en teclado, con algunos apoyos adicionales para
ratón en menús e interfaz. Los controles más importantes son los siguientes:

\begin{itemize}
    \item \textbf{Movimiento:} \texttt{W}, \texttt{A}, \texttt{S}, \texttt{D}.
    \item \textbf{Ataque:} flechas de dirección. En la configuración actual el combate queda orientado principalmente a este esquema de entrada.
    \item \textbf{Interacción:} \texttt{E}.
    \item \textbf{Cambio de forma:} \texttt{Tab}.
    \item \textbf{Pausa:} \texttt{Esc}.
    \item \textbf{Selección rápida de arma:} \texttt{1}, \texttt{2} y \texttt{3}.
    \item \textbf{Seleccionar siguiente hueco del inventario:} \texttt{Alt}.
    \item \textbf{Soltar objeto del hueco seleccionado:} \texttt{Q}.
    \item \textbf{Recuperar vida:} \texttt{Ctrl}, si el jugador no está a vida completa y dispone de almas suficientes.
    \item \textbf{Dash:} \texttt{Shift}, siempre que el jugador posea el objeto especial correspondiente.
\end{itemize}

\section{Interfaz de usuario}
Durante la partida la interfaz muestra el estado principal del jugador y del progreso:

\begin{itemize}
    \item La \textbf{barra de vida} permite seguir el estado actual del personaje.
    \item El \textbf{contador de almas} indica el recurso acumulado, utilizado para compras y ciertas acciones especiales.
    \item El \textbf{indicador de círculo} informa del nivel temático en el que se encuentra la \textit{run}.
    \item El \textbf{inventario} muestra los huecos disponibles, el hueco activo y los objetos almacenados.
    \item El \textbf{minimapa} representa la exploración de la planta actual y diferencia entre salas ocultas, visitadas y sala activa.
    \item El \textbf{indicador de llave} informa de si la llave del círculo actual ya ha sido recogida.
\end{itemize}

\section{Interacción con el mundo}
\subsection{Exploración y combate}
Al entrar en determinadas salas se activa un encuentro. Mientras haya enemigos vivos, las
puertas permanecen cerradas. Cuando el encuentro termina, la sala vuelve a abrirse y el
jugador puede continuar avanzando.

El personaje puede alternar entre forma de melé y forma a distancia. Esta dualidad afecta al
tipo de arma activa, al comportamiento ofensivo y a la lectura táctica de los encuentros.

\subsection{Objetos, tienda y tesoro}
Los objetos pueden obtenerse de varias formas:

\begin{itemize}
    \item Recogiéndolos del suelo al interactuar con \texttt{E}.
    \item Tomándolos de una sala de tesoro con \texttt{E}, siempre que exista espacio disponible en el inventario.
    \item Comprándolos en la tienda con \texttt{E}, siempre que el jugador disponga de suficientes almas y de un hueco libre.
\end{itemize}

El inventario utiliza un número limitado de huecos desbloqueados. Si está lleno, el juego no
permitirá recoger o comprar nuevos objetos hasta liberar espacio o aumentar la capacidad.

\subsection{Llave de círculo y descenso}
Cada círculo incorpora una llave asociada. Para progresar correctamente, el jugador debe
recoger la llave del círculo actual y después alcanzar la escalera. Mientras la llave no se
haya obtenido, la escalera permanece desactivada. Una vez recogida, queda habilitada y
permite descender al siguiente círculo al entrar en contacto con ella.

\subsection{Fuente bautismal}
La fuente bautismal permite transferir almas al jugador al interactuar con \texttt{E} dentro
de su radio de acción. Su función es ofrecer un punto de recuperación de recursos dentro de
la exploración, siempre que todavía disponga de almas almacenadas.

\section{Menús de pausa, victoria y derrota}
\subsection{Menú de pausa}
Durante la partida puede abrirse el menú de pausa con \texttt{Esc} o con el botón
\textit{Start} del mando. Al activarlo, el juego detiene el tiempo, pausa el audio y oculta
la interfaz principal de juego. Desde este menú pueden utilizarse tres botones:

\begin{itemize}
    \item \textbf{Reanudar:} vuelve a la partida actual.
    \item \textbf{Reiniciar:} reinicia la \textit{run} desde el principio.
    \item \textbf{Salir:} vuelve al menú principal.
\end{itemize}

\subsection{Pantalla de victoria}
Cuando se completa el último círculo, el juego muestra una pantalla de victoria. Desde ella
puede comenzarse una nueva partida mediante el botón de reinicio o volver al menú principal
mediante el botón de salida.

\subsection{Pantalla de derrota}
Si la vida del jugador llega a cero, aparece la pantalla de derrota. Sus opciones son
análogas a las de victoria:

\begin{itemize}
    \item \textbf{Reiniciar:} inicia una nueva \textit{run}.
    \item \textbf{Salir:} vuelve al menú principal.
\end{itemize}

\section{Modo de desarrollo}
Si la partida se inicia desde el botón \textbf{Dev Mode}, el juego activa un conjunto de
atajos adicionales orientados a pruebas y depuración. Estas acciones no forman parte del uso
normal previsto para un jugador final, pero resultan útiles durante demostraciones internas y
validación de sistemas.

\subsection{Atajos de depuración generales}
\begin{itemize}
    \item \textbf{\texttt{Z}:} añade almas al jugador.
    \item \textbf{\texttt{X}:} resta almas al jugador.
    \item \textbf{\texttt{Ctrl + N}:} fuerza el descenso al siguiente círculo.
\end{itemize}

\subsection{Atajos de objetos de depuración}
Si la configuración de la partida mantiene activa la inyección de objetos de prueba, el modo
de desarrollo permite generar rápidamente varios objetos especiales:

\begin{itemize}
    \item \textbf{\texttt{F1}:} objeto especial de \textit{Dash}.
    \item \textbf{\texttt{F2}:} objeto especial de \textit{Cloak}.
    \item \textbf{\texttt{F3}:} objeto especial de \textit{Shield}.
    \item \textbf{\texttt{F4}:} objeto especial de revelado completo del mapa.
    \item \textbf{\texttt{F5}:} objeto especial \textit{Cursed Coin}.
    \item \textbf{\texttt{F6}:} añade el siguiente objeto del \textit{pool} de depuración configurado.
\end{itemize}

\subsection{Atajos de tienda y tesoro}
Dentro de salas de tienda o tesoro, el modo de desarrollo activa además dos accesos rápidos:

\begin{itemize}
    \item \textbf{\texttt{F7}:} repuebla la tienda o el tesoro de la sala actual.
    \item \textbf{\texttt{F8}:} vacía el contenido de la tienda o del tesoro actual.
\end{itemize}

